\section{สรุปและอภิปราย}
คลังข้อมูลการลานิสิตใช้ข้อเท็จจริงระดับใบลาและระดับใบลากับกลุ่มเรียนร่วมกับมิติกลาง เพื่อรองรับคำถามตั้งแต่พฤติกรรมรายบุคคลถึงภาพรวมระดับคณะ จุดสำคัญของการออกแบบคือการกำหนด grain ให้ชัด แยกบริบทองค์กร และควบคุมการนับซ้ำจากรายวิชาหรือผู้สอนร่วม

มุมมองทั้งหกบทบาทใช้ข้อมูลพื้นฐานร่วมกัน แต่ต่างกันที่ขอบเขตประชากร รายละเอียดที่เปิดดูได้ และวัตถุประสงค์ การใช้ข้อมูลเพื่อการตัดสินใจจึงต้องระบุหน่วยนับ ช่วงเวลา ตัวกรอง และแหล่งข้อมูลกำกับทุกครั้ง

งานที่ควรดำเนินการต่อคือเชื่อมหน้าเว็บกับบริการจริง ตรวจยอดระดับใบลาและรายวิชาเทียบต้นทาง รักษาค่าไม่ทราบให้ต่างจากศูนย์ ตรวจสิทธิ์การเปิดรายละเอียดข้ามบุคคลและข้ามหน่วยงาน รวมถึงทดสอบการอ่านข้อมูลระหว่างเปลี่ยนรอบ ETL ทั้งนี้การจัดทำบทความครั้งนี้ไม่ได้รัน ETL โหลดฐานข้อมูล หรือวัดประสิทธิภาพระบบ จึงไม่มีตัวเลขผลทดลองหรือผลประเมินความพึงพอใจที่ยืนยันได้
